% ============================================================================
% Fichier  : partie2/P02_C05_analyse_memoire.tex
% Titre    : L'Analyse de la Mémoire
% Auteur   : MINKA MI NGUIDJOÏ Thierry Emmanuel
% Version  : 1.0 -- novembre 2026
% Corrigés : corriges/Corr_C05.tex
% ============================================================================
% Sources : Ligh et al. (2014) ; Volatility 3 (greffons windows.info,
% pslist, psscan, pstree, cmdline, netscan, malfind, svcscan) ; structures
% EPROCESS et VAD de Windows (présentées sans citation de la documentation).
% Faits nouveaux (bible) : sur SRV-AD01 (capture du 14/03 à 10 h 02),
% processus « svchost.exe » PID 4412 lancé depuis C:\ProgramData, parent
% services.exe, service WinUpdSvc créé le 09/03 à 1 h 35 ; connexion établie
% vers 203.0.113.47:443 ; fichier lsass.dmp créé le 09/03 à 1 h 37 sur le
% disque de SRV-AD01 par rundll32.exe (commande retrouvée dans les journaux).
% ============================================================================

\chapter{L'Analyse de la Mémoire}
\label{chap:memoire}
\chaptermark{L'analyse de la mémoire}

\epigraph{%
    « Un programme peut effacer son fichier.\\
    Il ne peut pas tourner sans être en mémoire. »%
}{Principe du cours}

\niveauChapitre{Intermédiaire}{%
    Maîtrise des Chapitres~\ref{chap:acquisition} et~\ref{chap:disque}
}

\begin{boxobjectifs}
\begin{itemize}
    \item Expliquer ce que la mémoire vive contient et que le disque ne
          garde pas.
    \item Identifier le profil d'une capture et en vérifier la cohérence.
    \item Lister les processus par deux méthodes et interpréter leurs
          écarts.
    \item Repérer un processus anormal par son parent, son chemin et sa
          ligne de commande.
    \item Relier un processus à ses connexions réseau et à sa persistance.
    \item Reconnaître les traces d'un vol d'identifiants, sans en
          reproduire la technique.
\end{itemize}
\end{boxobjectifs}

\bigskip

% ============================================================================
\section*{Avant de commencer : un svchost de trop}
% ============================================================================

Yaoundé, 22 mars 2025. M.~ABOUBAKAR, l'expert réquisitionné par le
procureur, a reçu le scellé NC-01, l'a ouvert devant témoin, a vérifié son
empreinte et travaille sur sa copie~: 16~Go de mémoire du contrôleur de
domaine, capturés le 14~mars à 10~h~02. Il compte, dans la liste des
processus, quatorze \texttt{svchost.exe}. Treize sont lancés depuis
\texttt{C:\textbackslash Windows\textbackslash System32}. Le quatorzième
tourne depuis \texttt{C:\textbackslash ProgramData}, et il est le seul à
entretenir une connexion vers l'extérieur.

\begin{boxalert}
\textbf{Violation typique.} Conclure qu'un processus est légitime parce que
son nom est celui d'un programme du système.

\smallskip
\textbf{Règle~:} un processus se juge sur son chemin, son parent, sa ligne de
commande, son heure de démarrage et son comportement réseau, comparés à ce
qu'on attend de ce programme. Le nom est le premier attribut qu'un attaquant
imite.
\end{boxalert}

\separateur

% ============================================================================
\section{Ce que la Mémoire Contient}
\label{sec:c05-contenu}
% ============================================================================

La mémoire vive d'un serveur Windows contient, à l'instant de la capture,
l'état complet du système~: les processus et leurs fils d'exécution, les
bibliothèques chargées, les connexions réseau, les objets ouverts (fichiers,
clés de registre), le cache des fichiers, une partie du registre, et les
données que les programmes manipulent, y compris des mots de passe, des clés
de chiffrement et des commandes saisies. Certains programmes malveillants
n'existent qu'en mémoire~: ils sont chargés depuis le réseau, ou injectés
dans un processus légitime, sans jamais être écrits sur le disque.

\begin{tableaucours}{L{4.4cm}L{3.6cm}L{4.2cm}}{%
    \tableauHeader{Question} &
    \tableauHeader{Mémoire} &
    \tableauHeader{Disque}
}
Quels programmes tournaient~? & Oui, avec leurs parents & Indirectement
(Amcache, journaux) \\
Quelles connexions étaient ouvertes~? & Oui, avec le processus & Non \\
Quelle ligne de commande~? & Oui & Seulement si l'audit l'enregistre \\
Code injecté dans un processus~? & Oui & Rarement \\
Ce qui s'est passé il y a un mois~? & Non & Oui, par les journaux et
les dates \\
\end{tableaucours}
\vspace{-0.8em}
\begin{center}{\small\itshape Tableau~5.1. Mémoire et disque~: deux sources complémentaires}\end{center}

% ============================================================================
\section{L'Outil et la Méthode}
\label{sec:c05-outil}
% ============================================================================

L'outil de référence, libre, est Volatility, dans sa version~3~\cite{volatility3}.
Il reconstruit les structures du système d'exploitation à partir de la
capture brute, grâce à des tables de symboles propres à chaque version du
noyau. On l'emploie greffon par greffon, chacun répondant à une question.

\begin{tableaucours}{L{3.6cm}L{8.6cm}}{%
    \tableauHeader{Greffon} &
    \tableauHeader{Question}
}
\texttt{windows.info} & La capture est-elle lisible, de quelle version de
Windows, à quelle heure système~? \\
\texttt{windows.pslist} & Quels processus figurent dans la liste tenue par
le noyau~? \\
\texttt{windows.psscan} & Quelles structures de processus trouve-t-on en
balayant toute la mémoire, y compris des processus terminés ou cachés~? \\
\texttt{windows.pstree} & Qui a lancé qui~? \\
\texttt{windows.cmdline} & Avec quelle ligne de commande~? \\
\texttt{windows.netscan} & Quelles connexions et quels ports, rattachés à
quel processus~? \\
\texttt{windows.malfind} & Quelles zones de mémoire exécutables et
modifiables, sans fichier associé, suggèrent une injection de code~? \\
\texttt{windows.svcscan} & Quels services sont enregistrés, avec quel
binaire~? \\
\end{tableaucours}
\vspace{-0.8em}
\begin{center}{\small\itshape Tableau~5.2. Les greffons de base de Volatility~3}\end{center}

\subsection{Deux listes de processus}

Windows tient la liste des processus actifs dans une liste chaînée de
structures \texttt{EPROCESS}. Le greffon \texttt{pslist} suit cette liste.
Le greffon \texttt{psscan}, lui, cherche ces structures partout dans la
capture, par leur forme. La comparaison des deux est instructive~: un
processus présent dans \texttt{psscan} et absent de \texttt{pslist} est soit
terminé (sa structure traîne encore en mémoire, avec une heure de fin), soit
caché par un programme qui l'a retiré de la liste. Le premier cas est
banal~; le second est un signe fort de compromission.

\subsection{Le processus normal}

Pour juger un processus, on compare ses attributs à ce que l'on sait des
processus du système. Quelques repères suffisent pour les plus imités.

\begin{tableaucours}{L{2.6cm}L{3.8cm}L{2.6cm}L{3.2cm}}{%
    \tableauHeader{Nom} &
    \tableauHeader{Chemin attendu} &
    \tableauHeader{Parent attendu} &
    \tableauHeader{Nombre}
}
\texttt{lsass.exe} & \texttt{System32} & \texttt{wininit.exe} & Un seul \\
\texttt{services.exe} & \texttt{System32} & \texttt{wininit.exe} & Un seul \\
\texttt{svchost.exe} & \texttt{System32} & \texttt{services.exe} &
Plusieurs \\
\texttt{explorer.exe} & \texttt{Windows} & Variable (lancé à la session) &
Un par session \\
\end{tableaucours}
\vspace{-0.8em}
\begin{center}{\small\itshape Tableau~5.3. Quelques processus et leurs attributs attendus}\end{center}

% ============================================================================
\section{Ce que la Mémoire de SRV-AD01 Révèle}
\label{sec:c05-nyanga}
% ============================================================================

\begin{boxscenario}{Opération NYANGA. Le processus 4412}
\textbf{Constatations} (copie de NC-01, Volatility~3).
\begin{itemize}
    \item \texttt{windows.info}~: Windows Server 2019, heure système de la
          capture 14/03/2025 09:02:11 UTC, soit 10:02:11 heure locale,
          cohérente avec le registre des actions.
    \item \texttt{pslist} et \texttt{psscan}~: quatorze processus
          \texttt{svchost.exe}~; le PID~4412 a pour chemin
          \texttt{C:\textbackslash ProgramData\textbackslash
          WinUpd\textbackslash svchost.exe}, pour parent
          \texttt{services.exe}, et a démarré le 09/03 à 1~h~35.
    \item \texttt{svcscan}~: un service \texttt{WinUpdSvc}, en démarrage
          automatique, a pour binaire ce même fichier.
    \item \texttt{netscan}~: le PID~4412 maintient une connexion établie vers
          203.0.113.47, port~443.
    \item \texttt{malfind}~: aucune zone suspecte dans \texttt{lsass.exe}~;
          deux zones exécutables sans fichier associé dans le PID~4412.
\end{itemize}

\textbf{Analyse} (confiance élevée). Le PID~4412 est un programme étranger
au système, déguisé sous le nom \texttt{svchost.exe}, installé comme service
pour redémarrer avec la machine, et en communication avec l'adresse vers
laquelle sont partis les 3,2~Go. C'est l'outil d'accès à distance de
l'attaquant sur le contrôleur de domaine.
\end{boxscenario}

Le parent \texttt{services.exe} ne disculpe pas le processus~: un service
est toujours lancé par \texttt{services.exe}, y compris un service créé par
un attaquant. C'est le chemin qui trahit le PID~4412, et le service qui
explique son parent.

\begin{boxobjet}{Processus 4412, capture mémoire de SRV-AD01}
\textbf{Source~:} scellé NC-01, copie de travail, Volatility~3.

\textbf{Attributs~:} nom \texttt{svchost.exe}~; chemin
\texttt{C:\textbackslash ProgramData\textbackslash WinUpd}~; parent
\texttt{services.exe}~; démarrage 09/03/2025 1~h~35~; service
\texttt{WinUpdSvc}~; connexion vers 203.0.113.47:443.

\textbf{Ce qu'il prouve~:} la présence, le 14~mars, d'un programme
persistant étranger au système, en communication avec l'adresse de
l'exfiltration.

\textbf{Ce qu'il ne prouve pas~:} comment le programme est arrivé sur la
machine, ni qui le pilotait.
\end{boxobjet}

% ============================================================================
\section{Le Vol d'Identifiants}
\label{sec:c05-identifiants}
% ============================================================================

Sous Windows, le processus \texttt{lsass.exe} gère les authentifications et
garde en mémoire des éléments d'identification des sessions ouvertes. Un
attaquant qui obtient les droits d'administration sur une machine peut
copier la mémoire de ce processus dans un fichier, puis l'emporter et
l'analyser chez lui pour en extraire des identifiants. Le cours ne décrit
pas la technique d'extraction~: elle relève de l'UE MSN-2S3-SEC03. Il
apprend à en reconnaître les traces.

\begin{tableaucours}{L{4.4cm}L{7.8cm}}{%
    \tableauHeader{Trace} &
    \tableauHeader{Où la chercher}
}
Fichier de vidage (\texttt{.dmp}) de \texttt{lsass.exe} & Disque~:
\texttt{\$UsnJrnl}, MFT, dossiers temporaires \\
Ligne de commande d'un outil de vidage & Mémoire (\texttt{cmdline}) si le
processus tourne encore~; événement 4688 si l'audit est actif \\
Accès au processus \texttt{lsass.exe} par un autre processus & Journaux de
surveillance avancée, s'ils existent \\
Réutilisation des identifiants & Ouvertures de session du compte volé depuis
des machines inhabituelles (4624) \\
\end{tableaucours}
\vspace{-0.8em}
\begin{center}{\small\itshape Tableau~5.4. Les traces d'un vol d'identifiants}\end{center}

\begin{boxscenario}{Opération NYANGA. Le fichier du 9 mars}
\textbf{Constatations.} Sur l'image du disque de SRV-AD01 (acquise le
16~mars), le \texttt{\$UsnJrnl} enregistre la création, le 9~mars à 1~h~37,
d'un fichier \texttt{lsass.dmp} de 54~Mo dans
\texttt{C:\textbackslash Windows\textbackslash Temp}, et sa suppression à
1~h~49. Le journal système enregistre, à 1~h~35, la création du service
\texttt{WinUpdSvc} (événement 7045) par le compte \texttt{svc\_backup}. À
partir du 10~mars, le compte \texttt{adm.kamga} ouvre des sessions depuis
SRV-AD01 vers les autres serveurs.

\textbf{Analyse} (confiance élevée). L'attaquant, entré sous
\texttt{svc\_backup}, a installé son outil le 9~mars, copié la mémoire de
\texttt{lsass.exe} deux minutes plus tard, puis utilisé les identifiants de
\texttt{adm.kamga}, présents en mémoire parce que l'administrateur s'était
connecté au contrôleur de domaine dans les jours précédents. Cette séquence
explique l'usage du compte sans la participation de son titulaire~; elle ne
l'exclut pas.
\end{boxscenario}

\begin{boxlecture}{La mémoire de SRV-AD01 et le disque associé}
\begin{tableaucours}{L{7.2cm}c}{%
    \tableauHeader{Proposition} &
    \tableauHeader{Cote}
}
Un programme étranger tournait sur SRV-AD01 le 14~mars & \OUI{} \\
Il communiquait avec 203.0.113.47 & \OUI{} \\
Il était installé comme service depuis le 9~mars & \OUI{} \\
La mémoire de \texttt{lsass.exe} a été copiée le 9~mars & \OUI{} \\
Les identifiants de \texttt{adm.kamga} ont été extraits de cette copie &
\NEK{} \\
Du code a été injecté dans \texttt{lsass.exe} & \NON{} \\
\end{tableaucours}

\medskip
\textbf{Lecture.} Quatre \OUI{}, un \NON{}, un \NEK{}. L'extraction des
identifiants depuis la copie reste \NEK{}~: elle a eu lieu hors du
système, chez l'attaquant, et aucune pièce ne peut la montrer. On l'infère
de la séquence, avec une confiance élevée, mais on ne la constate pas.
\end{boxlecture}

% ============================================================================
\section{Les Limites de l'Analyse Mémoire}
\label{sec:c05-limites}
% ============================================================================

\begin{itemize}
    \item Une capture est un instantané~: elle ne dit rien de ce qui a
          tourné puis s'est terminé, sauf traces résiduelles.
    \item Les tables de symboles doivent correspondre exactement à la
          version du noyau~; une version inconnue de l'outil donne des
          résultats faux ou vides.
    \item La capture elle-même peut être incohérente, parce que la mémoire a
          changé pendant qu'on la lisait~: des structures à moitié copiées
          apparaissent.
    \item Les zones signalées par \texttt{malfind} ne sont pas toutes
          malveillantes~: certains programmes légitimes (machines
          virtuelles Java, navigateurs) génèrent du code en mémoire.
\end{itemize}

% ============================================================================
\section{Ce que l'Analyse Mémoire demande à chaque Métier}
\label{sec:c05-metiers}
% ============================================================================

\begin{boxmetier}{Analyste en réponse à incident}
\textbf{Examiner une capture mémoire}
\begin{enumerate}
    \item Vérifier le profil et l'heure système avec \texttt{windows.info}.
    \item Comparer \texttt{pslist} et \texttt{psscan}.
    \item Examiner l'arbre des processus, les chemins et les lignes de
          commande.
    \item Rattacher chaque connexion à son processus.
    \item Chercher la persistance (services, tâches) du processus suspect.
    \item Confirmer chaque découverte par une autre source (disque, réseau).
\end{enumerate}
\end{boxmetier}

\begin{boxmetier}{Application OPJ/Expert Judiciaire}
L'expert indique la version de l'outil et des tables de symboles, la liste
des greffons et leur ordre, et conserve leurs sorties brutes en annexe~: un
autre expert doit pouvoir refaire l'examen sur la même copie et obtenir les
mêmes listes.
\end{boxmetier}

\begin{boxmetier}{RSSI}
Limiter les connexions des comptes d'administration du domaine aux seules
machines d'administration~: un identifiant qui ne se trouve pas en mémoire
d'un serveur ne peut pas y être volé.
\end{boxmetier}

\separateur

\begin{boxresume}
\textbf{Ce qu'il faut retenir de ce chapitre~:}
\begin{itemize}
    \item La mémoire montre ce qui tournait, avec quels parents, quelles
          commandes et quelles connexions~; le disque montre l'histoire.
    \item Volatility s'emploie greffon par greffon, chacun répondant à une
          question~; \texttt{pslist} et \texttt{psscan} se comparent.
    \item Un processus se juge sur son chemin, son parent, sa commande, son
          heure et son réseau, jamais sur son nom.
    \item Le vol d'identifiants se reconnaît à ses traces (fichier de
          vidage, réutilisation du compte)~; l'extraction elle-même reste
          souvent \NEK{}.
    \item À NYANGA, le PID~4412, faux \texttt{svchost.exe} installé comme
          service le 9~mars, communique avec 203.0.113.47.
\end{itemize}
\end{boxresume}

\begin{boxinfo}
\textbf{Points de contrôle.} Vous avez compris ce chapitre si vous savez~:
\begin{itemize}
    \item interpréter un écart entre \texttt{pslist} et \texttt{psscan}~;
    \item dire pourquoi un processus est suspect à partir de ses attributs~;
    \item distinguer ce qui est constaté de ce qui est inféré dans un vol
          d'identifiants.
\end{itemize}
\end{boxinfo}

% ============================================================================
\section*{Exercices}
% ============================================================================

\begin{boxexo}[title={\ding{45}\quad Exercice~5.1. Normal ou suspect \niveaubase}]
Pour chaque processus, dites s'il est normal ou suspect, et pourquoi~:
(a) \texttt{lsass.exe}, \texttt{System32}, parent \texttt{wininit.exe}~;
(b) \texttt{lsass.exe}, \texttt{System32}, parent \texttt{explorer.exe}~;
(c) \texttt{svchost.exe}, \texttt{System32}, parent \texttt{services.exe}~;
(d) \texttt{svchost.exe}, \texttt{Users\textbackslash Public}, parent
\texttt{cmd.exe}~; (e) deux processus \texttt{lsass.exe}.
\end{boxexo}

\begin{boxexo}[title={\ding{45}\quad Exercice~5.2. pslist et psscan \niveaumoyen}]
\texttt{psscan} trouve trois processus absents de \texttt{pslist}~: un
\texttt{rundll32.exe} terminé le 9~mars à 1~h~38~; un \texttt{cmd.exe}
terminé le 13~mars~; un \texttt{svch0st.exe} sans heure de fin.
\begin{enumerate}[(a)]
    \item Interprétez chacun.
    \item Lequel demande une vérification immédiate, et laquelle~?
    \item Que pourrait apporter la ligne de commande du \texttt{rundll32.exe},
          si elle est encore lisible~?
\end{enumerate}
\end{boxexo}

\begin{boxexo}[title={\ding{45}\quad Exercice~5.3. Relier les sources \niveaumoyen}]
Construisez un tableau qui relie, pour le PID~4412, chaque attribut constaté
en mémoire à une confirmation sur le disque de SRV-AD01 ou dans la capture
réseau. Indiquez pour chaque ligne si la confirmation existe, manque, ou
reste à chercher.
\end{boxexo}

\begin{boxexo}[title={\ding{45}\quad Exercice~5.4. Ce qu'on ne peut pas montrer \niveauexpert}]
L'avocat de M.~KAMGA soutient que la copie de la mémoire de
\texttt{lsass.exe} « ne prouve pas » le vol de ses identifiants.
\begin{enumerate}[(a)]
    \item A-t-il raison sur le plan de la constatation~? Sur celui de
          l'analyse~?
    \item Rédigez le passage du rapport d'expertise qui traite ce point.
    \item Quelle mesure, en place avant l'incident, aurait permis de
          constater l'usage des identifiants volés plutôt que de l'inférer~?
\end{enumerate}
\end{boxexo}

\bigskip

\begin{boxinfo}
\textbf{Pour aller plus loin.}
\begin{itemize}
    \item M.~H.~Ligh et al., \textit{The Art of Memory
          Forensics}~\cite{ligh2014memory}.
    \item La documentation de Volatility~3~\cite{volatility3}.
\end{itemize}
\end{boxinfo}

\begin{boxinfo}
\textbf{Transition.} Le PID~4412 parle à 203.0.113.47. Qu'a-t-il envoyé, et
quand~? La capture de la sonde réseau le dit~: Chapitre~\ref{chap:reseau}.
\end{boxinfo}
